\section{Evaluation}\label{sec:tsps-evaluation}

To assess the concrete performance of our protocols, we implement and benchmark
the schemes of \Cref{sec:tsps-abe} and \Cref{sec:tsps-kiltz}.

\paragraph{Realising $\algo{MPC}$.} \Cref{sec:tsps-mpc} describes the interfaces
our constructions assume for secure MPC over $\Zp$ and over a generic group
$\group{}$. Those over $\Zp$ are provided directly by existing MPC frameworks.
Those over $\group{}$ are not. Prior work adds elliptic curve group operations
to MP-SPDZ and shares group elements between two
parties~\cite{ESORICS:DOKSS20,NDSS:GPVRNC23}. What our schemes additionally
require, and what we implement, is a multiparty sharing of group elements
together with the mixed-domain multiplication $\algo{ExpGrp}$ described in
\Cref{sec:tsps-mpc}, including its authentication components in the malicious
setting. We benchmark two concrete frameworks built on this extension: ATLAS, in
the honest majority setting, and MASCOT, in the dishonest majority
setting~\cite{C:GLOPS21,CCS:KelOrsSch16}.

\paragraph{ATLAS.} ATLAS is a secret sharing based MPC scheme designed for the
honest majority setting, achieving unconditional security against both honest
but curious and malicious adversaries~\cite{C:GLOPS21}. It builds on an existing
protocol, improving its round complexity by a factor of two by combining ideas
from Beaver triples with existing multiplication techniques.  Parties can then
use one round of communication to cover two rounds' worth of multiplication
gates~\cite{C:DamNie07}. In terms of \Cref{sec:tsps-mpc}, linearity is preserved
and addition is therefore local, since ATLAS is built from a linear secret
sharing scheme. Multiplication is performed by consuming correlated randomness
prepared during pre-processing. Shares are combined by broadcasting them to all
parties. Because the underlying sharing is Shamir, it gives genuine
$(t+1)$-out-of-$n$ reconstruction for any $t < n/2$, and hence robustness in the
sense of \Cref{sec:tsps-construction}.

\paragraph{MASCOT.} MASCOT performs authenticated secure multiplication of
shared field elements in the maliciously secure dishonest majority setting. It
generates the multiplication triples this relies on using oblivious transfer,
extended to the arithmetic setting, together with MACs and standard sacrifice
techniques, so that the protocol either correctly outputs the product of the
shared elements or aborts once malicious behaviour is
detected~\cite{CCS:KelOrsSch16}. In terms of \Cref{sec:tsps-mpc}, addition is
local, since MASCOT is also built over a linear secret sharing scheme.
Multiplication consumes the triples generated via oblivious transfer.  Shares
are combined by broadcasting them, with MACs checked to ensure the
reconstruction is correct. Its sharing is additive between all $n$ parties,
which is the case $t = n-1$ of \Cref{sec:tsps-mpc}: the resulting threshold
signature is $n$-out-of-$n$, every signer must contribute to combining, and
robustness is replaced by abort. This is the trade the dishonest majority
setting makes, and it should be read alongside the timings below.

\paragraph{Implementation.} Our implementation builds on MP-SPDZ, a framework
for implementing MPC with pre-processing~\cite{CCS:Keller20}. MP-SPDZ is built
for the approach of sharing elements of $\Zp$, with all operations defined
between elements of $\Zp$. For the protocols of this chapter, we additionally
need operations between elements of $\group{1}$ and $\group{2}$.

We need the groups used to be pairing friendly. Although we do not distribute
any pairing operations, since these occur only in verification, we still need
all operations to take place over pairing friendly curves. The RELIC toolkit has
been integrated into MP-SPDZ for operations over pairing friendly curves and
shared computation over elements of the resulting elliptic curve
groups~\cite{NDSS:GPVRNC23,relictoolkit}.

% TODO: name the curve and the security level, and state the RELIC build
% configuration. Something like: "All benchmarks use the BLS12-381 curve at the
% 128-bit security level, with RELIC configured for ..." Reviewers will ask, and
% the signature and key sizes in \Cref{tab:tsps-comp} are only concrete once the
% curve is fixed.

MP-SPDZ operates with a pre-processing (offline) phase and an online phase. In
our setting, $\algo{OffSign}$ corresponds to the pre-processing phase, and
$\algo{OnSign}$ and $\TSPScombine$ run in the online phase.

We benchmark both the honest majority and dishonest majority settings, using
ATLAS and MASCOT respectively to realise multiplication of shared
secrets~\cite{C:GLOPS21,CCS:KelOrsSch16}. Both provide the interfaces of
\Cref{sec:tsps-mpc}, with $\algo{Inv}$ built from $\algo{Mul}$ as described
there, $\algo{RandShare}$ taking the role of \texttt{input}, and $\algo{Open}$
taking the role of \texttt{output}.

\paragraph{What the benchmarks do and do not show.} Three caveats should be read
alongside the numbers. First, ATLAS is designed to work in both the semi-honest
and malicious settings, but the implementation available in MP-SPDZ captures
only the semi-honest case. The honest majority timings are therefore for
semi-honest security, while the dishonest majority timings are for malicious
security~\cite{CCS:Keller20}. Second, neither realisation is proved secure
against adaptive corruption, so the concrete schemes achieve the static analogue
of the notion proved in \Cref{sec:tsps-security}. Third, MP-SPDZ constructs the
output within its online phase, so the reported $\algo{OnSign}$ timings include
the cost of $\TSPScombine$. The communication attributed to $\algo{OnSign}$ in
\Cref{tab:tsps-comms} is still zero, because $\algo{OnSign}$ itself is
local and all of the communication in that column belongs to combining.

We run benchmarks on a 16-core machine running RHEL 9.0-64 with 32GB of RAM,
inside Docker containers based on Ubuntu 22.

% TODO: state the network model explicitly. All parties are co-located on one
% host, so these figures are effectively a LAN lower bound with negligible
% latency. Either say so in one sentence, or add a simulated WAN row using tc
% netem (e.g. 50ms RTT, 100Mbps) -- round complexity is exactly where ATLAS and
% MASCOT differ, so a WAN column would strengthen the comparison considerably.

\begin{table}[t] \centering \resizebox{\textwidth}{!}{ \begin{tabular}{ c || c |
c | c || c | c | c } \toprule & Assumption & Message space & $|\sk|$ & $|\pk|$ &
$|\SPSsig|$ & Verification cost \\ \midrule ~\cite{AC:CKPSS23} & PS + ROM &
$(\gen^{\vec{\msg}}, \ggen^{\vec{\msg}}) \in \group{1}^\ell \times
\group{2}^\ell$ & $ (\ell+1) | \Zp | $ & $ \ell | \group{1} | + (\ell+1) |
\group{2} | $ & $ 2 | \group{1} | $ & $3\ell+2$ \\ ~\cite{PKC:MMSST24} & SXDH &
$ \vec{\msg} \in \group{1}^\ell$ & $ 2(\ell+1) | \Zp | $ & $ (\ell+1) |
\group{2} | $ & $ 6 | \group{1} | + | \group{2} | $ & $\ell+11$ \\
\Cref{sec:tsps-abe} & GGM & $ \vec{\msg} \in \group{1}^\ell$ & $ (\ell+1) | \Zp
| $ & $ | \group{1} | + \ell| \group{2} | $ & $ | \group{1} | + 2 | \group{2} |
$ & $\ell+4$ \\ \Cref{sec:tsps-kiltz} & SXDH & $ \vec{\msg} \in \group{1}^\ell$
& $ 2(\ell+1) |\Zp | + 6 | \group{1} | $ & $ (\ell+7)| \group{2} | $ & $ 6 |
\group{1} | + | \group{2} |$ & $\ell+9$ \\ \bottomrule \end{tabular} }
\caption{Comparison of threshold structure preserving signature schemes.
Verification cost lists the number of pairing evaluations needed during
verification.} \label{tab:tsps-comp} \end{table}

% TODO: recheck the two rows for prior work against their papers. The rows for
% our own schemes have been recomputed from the descriptions in
% \Cref{sec:tsps-instantiation}: for \Cref{sec:tsps-kiltz}, sk = (K, b, x, y)
% with K in Zp^{(l+1)x2} and b, x, y each in G_1^2, giving 2(l+1)|Zp| + 6|G_1|;
% pk = (a~, u~, v~, w~) with a~, v~, w~ in G_2^2 and u~ in G_2^{l+1}, giving
% (l+7)|G_2|; and verification is 2 + (l+1) + 2 + 2 pairings for the first
% equation and 2 for the second, giving l+9.

\begin{table}[t] \centering \begin{tabular}{ c || c | c | c } \toprule &
$\algo{OffSign}$ & $\algo{OnSign}$ & $\algo{Combine}$ \\ \midrule
~\cite{AC:CKPSS23} & -- & 0 & $ 2 | \group{1} |$ \\ ~\cite{PKC:MMSST24} & -- & 0
& $ 6 | \group{1} | + | \group{2} |$ \\ \Cref{sec:tsps-abe} & $ n(\ell+4)|\Zp|$
& 0 & $| \group{1} | + 2 | \group{2} |$ \\ \Cref{sec:tsps-kiltz} & $ n(12|\Zp| +
6|\group{1}|) $ & 0 & $6 | \group{1} | + | \group{2} | $ \\ \bottomrule
\end{tabular} \caption{Per party communication complexity for the signing and
combining phases of the threshold signing schemes. A dash indicates that the
scheme has no offline phase.} \label{tab:tsps-comms} \end{table}

% TODO: recheck the OffSign column against the corrected protocol boxes. For
% \Cref{sec:tsps-kiltz}, OffSign is now one Mul over Zp plus eight calls to
% ExpGrp_1 (two for each of b^rho, b^{rho sigma}, x^rho, y^{rho sigma}), where
% the previous draft said four; the entry n(12|Zp| + 6|G_1|) may need updating
% accordingly. For \Cref{sec:tsps-abe}, OffSign is one Inv and l+1 Muls over Zp
% with no group operations at all, which is worth stating in the text since it
% is why that scheme is cheaper offline.

\begin{table}[t] \centering \begin{tabular}{ c || c || c | c || c | c } \toprule
& Centralised & \multicolumn{2}{c||}{ATLAS~\cite{C:GLOPS21}} &
\multicolumn{2}c{MASCOT~\cite{CCS:KelOrsSch16}} \\ & & $\algo{OffSign}$ &
$\algo{OnSign}$ & $\algo{OffSign}$ & $\algo{OnSign}$ \\ \midrule
\Cref{sec:tsps-abe} & 10ms & 46ms & 56ms & 92ms & 84ms \\ \Cref{sec:tsps-kiltz}
& 40ms & 88ms & 61ms & 352ms & 94ms \\ \bottomrule \end{tabular}
\caption{Timings of the sub-protocols of our threshold structure preserving
signatures, averaged over 1000 runs, for $n = 5$ signers, corruption threshold
$t$, and message length $\ell$. The $\algo{OnSign}$ columns include the cost of
$\TSPScombine$.} \label{tab:tsps-benchmarks} \end{table}

% TODO: fill in the values of t and l used for \Cref{tab:tsps-benchmarks} (the
% caption currently names them without giving them), and add rows for the two
% prior schemes. The contribution paragraph of \Cref{sec:tsps-introduction}
% claims that we implement previously suggested threshold structure preserving
% signature schemes for the sake of concrete comparison, and at present no table
% reports timings for them, so either the rows or the claim has to go.

% TODO: two derived numbers would be worth adding here, both cheap to obtain
% from the existing runs. (i) Throughput: signatures per second with a warm
% pre-processing store, which is the headline figure for the bursty
% high-throughput applications the chapter motivates. (ii) Storage: bytes per
% stored pre-processing entry, which for \Cref{sec:tsps-abe} is (l+1)|Zp| +
% |G_1| + 2|G_2| per signer and for \Cref{sec:tsps-kiltz} is 7|G_1| + |G_2| per
% signer, so that a reader can size the store against an expected burst.

\Cref{tab:tsps-comp} compares our two schemes against two prior threshold
structure preserving signature schemes in terms of assumptions, message space,
key and signature size, and verification cost~\cite{AC:CKPSS23,PKC:MMSST24}.
\Cref{tab:tsps-comms} compares per party communication complexity across the
offline, online, and combining phases. \Cref{tab:tsps-benchmarks} reports
concrete timings for both of our schemes, using ATLAS and MASCOT to realise the
underlying MPC interfaces, and against a centralised (non-threshold) baseline,
for $n = 5$ parties.

\paragraph{What each scheme is for.} The tables invite an obvious question,
which we answer directly rather than leave to the reader. Both prior schemes
sign without any offline phase at all, and ours do not, so why would anyone use
ours? On communication we are strictly worse, and we begin there. A signature
from either prior scheme costs only the shares sent during combining, whereas
ours costs that plus an amortised share of the offline phase, and
\Cref{tab:tsps-comms} shows the difference. What is bought in exchange differs
by scheme.

Against the scheme of~\cite{AC:CKPSS23}, the difference is the message space.
Its signers require the message in both source groups, as an indexed pair
$(\gen^{\vec{\msg}}, \ggen^{\vec{\msg}})$, which means the signers can only sign
a message someone can supply in that form. Ours sign an arbitrary vector in
$\group{1}^\ell$, so they can sign a commitment or a public key that the signers
did not generate and whose discrete logarithms nobody knows, which is what the
applications of \Cref{sec:tsps-introduction} actually need.  Verification is
also markedly cheaper, at $\ell+4$ or $\ell+9$ pairings against $3\ell+2$, a gap
that widens with the message length.

Against the scheme of~\cite{PKC:MMSST24}, the two comparisons differ from each
other. The scheme of \Cref{sec:tsps-abe} is smaller on every axis: three group
elements against seven, a public key of $|\group{1}| + \ell|\group{2}|$ against
$(\ell+1)|\group{2}|$, and $\ell+4$ pairings to verify against $\ell+11$. Its
signature size is optimal for a type III pairing~\cite{C:AGHO11}. The price is
the assumption: a $q$-type assumption in the generic group model, against SXDH.
A reader who will accept the generic group model should use \Cref{sec:tsps-abe},
and one who will not should not.

The scheme of \Cref{sec:tsps-kiltz} is the harder case, and we do not claim an
efficiency advantage for it. It shares its assumption, message space and
signature size with the scheme of~\cite{PKC:MMSST24}, verifies slightly more
cheaply at $\ell+9$ pairings against $\ell+11$, and is worse on key size, at
$(\ell+7)|\group{2}|$ against $(\ell+1)|\group{2}|$, as well as requiring the
offline phase. The case for it is a different one. It is not a bespoke threshold
scheme but an instance of the generic construction of
\Cref{sec:tsps-construction} applied to an existing, well studied structure
preserving signature scheme, left otherwise untouched. By
\Cref{lem:tsps-transparency} its public keys, signatures and verification
equation are exactly those of Kiltz, Pan, and Wee. A deployment that already
verifies those signatures can therefore thresholdise its signer without changing
any verifier, and without asking anyone to accept a new scheme. The same
construction applies to any structure preserving signature admitting the
decomposition of \Cref{sec:tsps-construction}, so it is a route to
thresholdising future schemes rather than a single scheme. \Cref{sec:tsps-abe}
and \Cref{sec:tsps-kiltz} are two instances of it, and the fact that one is
competitive on size and one is not is a property of the underlying signatures
rather than of the construction.

Finally, the offline phase is not purely a cost. Both prior schemes do their
work at signing time, so their cost falls on the critical path of every
signature. Ours moves almost all of it ahead of the message being known, which
matters in exactly the setting \Cref{sec:tsps-introduction} describes: bursty
traffic, where pre-processing accumulated during a quiet period is spent during
a busy one. The comparison to make is therefore not offline cost against zero,
but latency at the moment a signature is needed, and there \Cref{tab:tsps-comms}
records an online communication cost of zero for all four schemes.

\paragraph{Comparing across total number of parties.} We next examine how
execution time varies with the total number of parties, fixing the threshold of
participating parties and the message length at two group elements. In each
case, ATLAS is around twice as fast as MASCOT, which is expected, since MASCOT
relies on public key cryptography to handle dishonest majorities.
\Cref{fig:eval-keygen} and \Cref{fig:eval-offline} show that ATLAS scales better
than MASCOT as the number of parties increases. Since our benchmarks build on
MP-SPDZ, they reflect the standard MPC offline and online phases, where the
output is constructed within the online phase. This means the online and
combining phases are benchmarked together, so their timings grow with the number
of parties. In practice, parties may instead perform only the constant-time
$\algo{OnSign}$ locally and submit their signature shares to a single party,
which performs $\TSPScombine$. This avoids every party communicating shares and
constructing its own copy of the signature. The signature generation phase shown
in \Cref{fig:eval-online} combines $\algo{OnSign}$ and $\TSPScombine$.

% TODO: fill

\begin{figure}[t] \centering \fbox{\parbox[c][6cm][c]{0.85\textwidth}{\centering
PLACEHOLDER: $\TSPSgen$ execution time against total number of parties $n$, for
ATLAS and MASCOT.}} \caption{Key generation time against the total number of
parties, for both realisations of $\algo{MPC}$.} \label{fig:eval-keygen}
\end{figure}

\begin{figure}[t] \centering \fbox{\parbox[c][6cm][c]{0.85\textwidth}{\centering
PLACEHOLDER: $\algo{OffSign}$ execution time against total number of parties
$n$, for ATLAS and MASCOT.}} \caption{Offline signing time against the total
number of parties, for both realisations of $\algo{MPC}$.}
\label{fig:eval-offline} \end{figure}

\begin{figure}[t] \centering \fbox{\parbox[c][6cm][c]{0.85\textwidth}{\centering
PLACEHOLDER: combined $\algo{OnSign}$ and $\TSPScombine$ execution time against
total number of parties $n$, for ATLAS and MASCOT.}} \caption{Signature
generation time, comprising $\algo{OnSign}$ and $\TSPScombine$, against the
total number of parties, for both realisations of $\algo{MPC}$.}
\label{fig:eval-online} \end{figure}
